\documentclass{article}

%maths
\usepackage{mathtools}
\usepackage{amsmath}
\usepackage{amssymb}
\usepackage{amsfonts}

%tikzpicture
\usepackage{tikz}
\usepackage{scalerel}
\usepackage{pict2e}
\usepackage{tkz-euclide}
\usetikzlibrary{calc}
\usetikzlibrary{patterns,arrows.meta}
\usetikzlibrary{shadows}

%pgfplots
\usepackage{pgfplots}
\pgfplotsset{compat=newest}
\usepgfplotslibrary{statistics}
\usepgfplotslibrary{fillbetween}

%colours
\usepackage{xcolor}

\usepackage{nicefrac}

\begin{document}

\section{Rational Expressions}
The quotient of two polynomials is called a \textbf{rational expression}, where the denominator is never zero.\\
\\
Examples include $5/6$, $\frac{3x + 8}{2x - 9}$, $6x + 5$\\
\\
The last one can be written as $\frac{6x + 5}{1}$ which fits the definition.

\subsection{Simplifying Rational Expressions}
1. Completely factor numerators and denominators.
2. Reduce common factors.\\
\\
\subsection{Multiplying Rational Expressions}
1. Completely factor all numerators and denominators.\\
2. Reduce all common factors.\\
3. Either multiply the denominators, using FOIL if needed, or leave it in factored form.\\
\\
\subsection{Dividing Rational Expressions}
Division of fractions involves multiplying by the \textit{reciprocal} of the divisor.

\subsection{Adding and Subtracting Rational Expressions}
If they have the same denominator then add/subtract as usual making sure to simplify at the end.\\
If they \textit{dont} have the same denominator then you need to find the Lowest Common Denominator (LCD).\\

\subsection{Complex Fractions}
If the numerator or denominator or both contain fractions the expression is called a \textbf{complex fraction.}\\
\\
1. Simplify both numerator and denominator expressions into single fraction expressions.\\
2. Divide the numerator by the denominator and simplify if possible.\\
\\

\subsection{Proportion, Direct Variation, Inverse Variation, Joint Variation}
\textbf{Proportion} is an equation stating that two rational expressions are equal. Simple ones are solved with the \textit{Cross Product Rule}\\
\\
\[
\fbox{
    \parbox{4cm}{ % You need to set a width for the parbox
        \textbf{Cross Product Rule} \\ 
        if $\frac{a}{b} = \frac{c}{d}$, then $ad = bc$.
    }
}
\]
\textbf{Direct Variation} is when 'y varies directly as x' or 'y is directly proportional to x'. This may be inteperted in two ways:\\
\\
1. $\frac{y}{x} = k $ \\
\\
For some constant k - k is called the \textbf{constant of proportionality}. This is used when the constant is the desired result.\\
\\
2. $\frac{y_1}{x_1} = \frac{y_2}{x_2}$ \\
\\
This is when the desired result is either an original or new value of x or y.\\
\\

\textbf{Inverse Variation} is when 'y varies inversely (opposite) to x'.

\textbf{Joint Variation} is when one variable varies as the product of the other variables. The phrase goes 'y varies jointly as x and z':\\
\\
1. $\frac{y}{xz} = k$ \\
\\
if the constant is desired.\\
\\
2. $\frac{y_1}{x_1z_1} = \frac{y_2}{x_2z_2}$\\
\\
if one of the variables is desired.\\
\\


\end{document}
